\section{Discussion}\label{sec:discusion}

\paragraph{What remains open for the exact count.}
Set $R_Q(n)=\log Q(n)-\log Q_{\rm pred}(n)$. Theorem~\ref{thm:intro-constante} is equivalent to
\[
 \log Q(n)-\log n!-n\Phi=C_{\rm pred}+R_Q(n)+o(1).
\]
Thus $R_Q(n)\to0$ suffices for the constant term of $\log Q(n)$ to equal $C_{\rm pred}$. If
all that were known was that $R_Q(n)$ has a finite limit $r$, the constant term would exist and would equal
$C_{\rm pred}+r$. We have proved neither statement. The natural quantitative form is the following.

\begin{conjetura}\label{conj:RQ}
There exist $C$ and $n_0$ such that $|\log Q(n)-\log Q_{\rm pred}(n)|\le C/n$ for all $n\ge n_0$.
\end{conjetura}

Section~\ref{sec:numerico} gives $R_Q(20)\approx0.021$ and $R_Q(27)\approx0.012$.

\paragraph{The role of consistency.}
Theorem~\ref{thm:intro-consistencia} states that, on residual boards in its classes, the predictor satisfies
the exact recursion for the count at each step up to a factor $e^{O(k^{-2})}$. Along a trajectory of
residual boards that
stays in a single class with $k\ge\delta_0n$, the errors sum to $O_{\delta_0}(1/n)$
(Section~\ref{sec:B4T05}). To pass from this to $R_Q(n)\to0$ requires three ingredients that this paper
does not provide: membership of typical residual boards of a random placement in these classes; control of
the end of the process, when few empty rows remain (for the exact count,~\cite{Completabilidad} studies the
case of
a fixed number of queens left to place, without uniformity as this number grows); and a comparison between
the count and the predictor along random trajectories.

The comparison is exact at each step. Let $R(B)=\log Q(B)-L(B)$, so that $R_Q(n)$ is its value on
the full board (Lemma~\ref{lem:puente}), and let $\tilde w_c=e^{L(B+c)-L(B)}/\rho(B)$, which are positive
weights summing to one over the children of row $r$. If $Q(B)>0$ and the predictor is defined at $B$
and at these children, the recursion for the count gives
\[
 R(B)=\log\rho(B)+\log\sum_{c\in r}\tilde w_c\,e^{R(B+c)},
\]
where a child without completions contributes zero to the sum. This identity indicates what would have
to be controlled to bound $R_Q(n)$: the values of $\log\rho$ along the tree of insertions, averaged
with these weights, and the value of $R$ on the residual boards where the predictor ceases to be
defined.

\paragraph{Finite activity.}
The predictor is defined directly on the full board, on the polytope $K_n$. At finite
activity $\lambda$, the corresponding approximation uses the small branch of the Bethe equations.
For $n\ge6$, existence of the small branch for every $\lambda>0$ is a conjecture. A floating-point
calculation, without rigorous bounds, finds it at the 61 activities tested, spaced
logarithmically between $\Lambda=\lambda(n-1)=10^{-3}$ and $10^{12}$, for $n=6$, $7$, $8$, $10$, $12$,
$16$, $20$, $27$, $40$, $60$, $80$, $120$ and $160$. For $n=4$, $K_4$ reduces to the point
$\pi\equiv1/4$, because each row has four squares with $p_c\le1/4$ and sums to one; hence
$(H_\infty)$ does not hold. For $n=5$, the calculation indicates $\max_c\pi^{(5)}_c=1/4$. In both cases, a
barrier optimizer suggests that the branch ceases to exist for large $\Lambda$, but this has not been
proved.
Nothing in this paragraph is used in the rest of the paper.
